% 18 November 2026. Drafted for Edward Baker with OpenAI Codex assistance.
% Actual model and effort unavailable; helper settings inherited.
\section{Canonical comparisons on a fixed positive-time slab}
\label{sec:canonicaltime}

We now acquire the preceding inputs from the unchanged canonical source.
There are three distinct conclusions. Complete initial interaction data
quantify the finite allowance of Theorem~\ref{stat:backward-comparison};
a growing head of evolving neighbors supplies the stronger joint barrier;
and spatial source jets give a sharper estimate of the actual gap drift.
All three use a fixed slab bounded away from zero. The published source
theorem already proves that its sufficiently high zeros are real and
simple throughout that slab. Consequently the conclusions refine the
dynamical comparisons on this known exterior, without extending its
coverage. We follow \cite{BakerCM21,BakerCM24,BakerCM25}, distinguishing
the published inputs recorded in \cite{BakerCMSource04} from their
moment and comparison consequences.

\subsection{The source clock and common labels}

Fix
\begin{equation}
 I=[t_0,t_*],\qquad 0<t_0\le t_*\le\tfrac12,
 \qquad \Delta=t_*-t_0.
 \label{ct:slab}
\end{equation}
The clock used below is
\begin{equation}
 G(x,t)=\frac{x}{4\pi}\log\frac{x}{4\pi}-\frac{x}{4\pi}
       +\frac{11}{8}+\frac{t}{16}\log\frac{x}{4\pi},
 \qquad G(q_n(t),t)=n.
 \label{ct:clock}
\end{equation}
Theorem~1.5(i)--(iv) of \cite{Polymath2019DBN}, with its absolute-constant
convention and the multiplicity count in its proof, supplies the following
inputs. For $0<t\le1/2$, sufficiently high zeros are exhausted by one
real simple occurrence $x_n(t)$ near each clock point, with
\begin{equation}
 x_n(t)=q_n(t)+O(q_n(t)^{-c_*t}).
 \label{ct:source-position}
\end{equation}
Here $c_*>0$ and the error constants are absolute; choose $c_*$ small
enough for both this localization and the source estimates used below.
The height and index
onsets are both exponential in $1/t$, and the theorem also supplies
the complete unit-window occurrence bound used in
\eqref{eq:count-safe}. We use the weaker displayed localization
\eqref{ct:source-position}; the source proof contains a stronger
displacement estimate. No numerical onset constant is extracted here.

\begin{lemma}[One cutoff and fixed moving labels]
\label{ct:labels}
After a single sufficiently large choice of onset depending on $t_0$,
the localization is uniform on $I$, the high labels are consecutive,
and each $x_n(t)$ is one $C^1$ branch throughout $I$. With spatial
buffers retained above that onset,
\begin{equation}
 q_n'(t)=-\frac{\pi}{4}
       \left(1+\frac{\pi t}{4q_n(t)\log(q_n(t)/(4\pi))}\right)^{-1},
 \qquad |q_n'(t)|\le\pi/4.
 \label{ct:clock-velocity}
\end{equation}
The same exterior statement holds at negative heights by evenness.
\end{lemma}

\begin{proof}
On $I$, the height threshold $\exp(C/t)$ is at most
$\exp(C/t_0)$ and $x^{-c_*t}\le x^{-c_*t_0}$. Increase the onset
to meet the index threshold as well: by \eqref{ct:clock}, an index at
height $x$ has size comparable to $x\log x$. This is an index-to-height
conversion, not an identification of the two thresholds. The same
increase pays any fixed buffer used below.

The source proof counts one occurrence in each sufficiently small disk
of radius $c/\log q_n$ and exhausts all zeros in the high region.
Clock separation is asymptotic to $4\pi/\log q_n$, so the disks are
disjoint after a further uniform enlargement. A count of one with
multiplicity proves simplicity; conjugation proves reality. Differentiating
the explicit equation $G(q_n(t),t)=n$, using
\[
 G_x(x,t)=\frac1{4\pi}\log\frac{x}{4\pi}+\frac{t}{16x},
 \qquad G_t(x,t)=\frac1{16}\log\frac{x}{4\pi},
\]
gives \eqref{ct:clock-velocity}. The moving disks are continuous and
their simple roots patch by the implicit-function theorem. Uniqueness
in each disk prevents a permutation of labels. Choosing the initial
label with the motion bound and the later window buffers included keeps
all disks in the source domain for the entire slab.
\end{proof}

This result does not assume reality of the low divisor. The later use
of positive exterior summands will explicitly impose the globally real
stopped regime. At $t_0=0$, the onset bound diverges and the error in
\eqref{ct:source-position} ceases to decay; the argument gives no
time-zero localization theorem, even though upper counts may remain
available there.

\subsection{Local geometry and the complete interaction}

Choose a buffered high label $n$, and freeze
\begin{equation}
 X=q_n(t_*),\qquad L=\log\frac{X}{4\pi}.
 \label{ct:scale}
\end{equation}
By Lemma~\ref{ct:labels}, $x_n(t)=X+O(1)$ on $I$. At a slice put
$a=x_n(t)$, $b=x_{n+1}(t)$, $g=b-a$, $D=g^2/4$, and
$d=1/G_x(a,t)$. In the next lemma take either $R=1$ or $R=L^2$, and set
\begin{equation}
 E_R=LX^{-c_*t_0}+\frac{R^2}{X}.
 \label{ct:window-error}
\end{equation}
All constants are independent of $X$ and the slice in the fixed slab.
They may depend on the stated source constants and slab. Every
asymptotic assertion is for $X$ sufficiently large; the threshold is
not evaluated numerically.

\begin{lemma}[Clock head and complete tail]
\label{ct:head-tail}
Uniformly for every actual high neighbor satisfying
$|x_{n+k}(t)-a|\le2R$,
\begin{equation}
 \frac{x_{n+k}(t)-a}{d}=k+e_k,
 \qquad |e_k|\le C E_R,
 \qquad e_0=0.
 \label{ct:normalized-position}
\end{equation}
In particular, the target and its two neighboring gaps are
$d(1+O(E_1))$, and $d\sim4\pi/L$.
At a slice where the complete divisor is real and simple, its complete
exterior interaction satisfies
\begin{equation}
 d^2M=2+O\left(E_R+\frac1{RL}\right),\qquad
 2DM=1+O\left(E_R+\frac1{RL}\right).
 \label{ct:complete-mass}
\end{equation}
Every exterior occurrence, including negative and bounded-height zeros,
is included in these estimates.
\end{lemma}

\begin{proof}
Throughout the radius-$2R$ window, $R=o(X)$ and
$G_x=O(L)$, $G_{xx}=O(1/X)$. Applying $G$ to
\eqref{ct:source-position} therefore gives
\[
 G(x_{n+k}(t),t)=n+k+O(LX^{-c_*t_0}).
\]
Subtract the corresponding equation at $a$. Taylor expansion about
$a$, with its full $O(R^2/X)$ curvature remainder, gives
\eqref{ct:normalized-position}. This uses values of the source error,
not derivatives of it. Consecutivity follows from source exhaustion;
subtracting neighboring normalized positions proves the gap assertions.

For the complete interaction, split the divisor at $|\rho-a|=R$,
assigning equality to the far set. The near indices form an interval
whose endpoints have magnitudes comparable to $RL$, with $0,1$ removed.
Indeed $G_x\asymp L$ on the window, and
\eqref{ct:normalized-position} has error $o(1)$. For such an exterior
index,
\begin{equation}
 \frac{d^2}{(\rho-a)(\rho-b)}
 =\frac1{(k+e_k)(k-1+e_k-e_1)}
 =\frac{1+O(E_R)}{k(k-1)}.
 \label{ct:head-summand}
\end{equation}
Both reference integer factors are nonzero and have magnitude at least
one, so the relative estimate is uniform over the whole head. The
positive reference sum and its two omitted tails are
\[
 \sum_{k\in\mathbb Z\setminus\{0,1\}}\frac1{k(k-1)}=2,
 \qquad
 \sum_{|k|>cRL}\frac1{|k(k-1)|}=O((RL)^{-1}).
\]
Consequently $d^2M_{\mathrm{near}}=2+O(E_R+(RL)^{-1})$.

For the far set let $r=|\rho-a|\ge R$. Eventually $g<1/4$, and
$|(\rho-a)(\rho-b)|\ge(3/4)r^2$. In each pair of unit shells
$R+j\le r<R+j+1$, $j\ge0$, the complete count bounds the number
of occurrences by $C\log(3+X+R+j)$. Hence
\begin{equation}
 |M_{\mathrm{far}}|
 \le C\sum_{j\ge0}
       \frac{\log(3+X+R+j)}{(R+j)^2}
 \le C'\frac{L+\log(R+1)+1}{R}
 =O(L/R).
 \label{ct:tail-payment}
\end{equation}
The middle inequality follows by splitting the logarithm into a
$\log(3+X)$ term and a $\log(1+R+j)$ term, and comparing the two
sums with the integrals of $r^{-2}$ and $(\log(1+r))r^{-2}$.
This accounts for every unbounded shell. If a shell crosses zero,
split its unit intervals and use reflection; the change is only in
the count constant. Boundary occurrences are retained by the stated
near/far convention.

Since $d^2=O(L^{-2})$, \eqref{ct:tail-payment} contributes
$O((RL)^{-1})$ to the normalized mass. This proves the first part
of \eqref{ct:complete-mass}; $g=d(1+O(E_R))$ and $2D=g^2/2$
give the second.
\end{proof}

The passage from local spacing to complete $M$ is essential. Local
clock data alone do not identify the complete interaction; the error
in \eqref{ct:tail-payment} pays the rest of the divisor. Conversely,
the positive lower-head argument for $K$ below will not claim an
asymptotic equality by silently discarding its tail.

\subsection{The exact initial-data allowance at high height}

At an all-real simple slice in $I$, let $h_L,h_R$ be the two neighboring
gaps and put
\[
 b_3=\frac1{h_L(h_L+g)}+\frac1{h_R(h_R+g)},\qquad
 \eta_X=LX^{-c_*t_0}+X^{-1},\qquad
 \varepsilon_X=L^{-1}+LX^{-c_*t_0}.
\]
The slice condition here is explicit: the following asymptotics are
uniform over all such slices, without presuming reality of the entire
slab.

\begin{theorem}[Initial-mass horizon]
\label{ct:initial-horizon}
For high pairs at the stated slices,
\begin{equation}
 g=\frac{4\pi}{L}(1+O(\eta_X)),\qquad
 DM=\frac12+O(\varepsilon_X),\qquad
 Db_3=\frac14+O(\eta_X).
 \label{ct:initial-asymptotics}
\end{equation}
Among all admissible initial lower data
$0<d_0\le D$ and $0\le b_0\le M$, the largest score in
\eqref{stat:strict-margin} is the exact-data score. Its strict duration
and the two-neighbor duration satisfy
\begin{equation}
 \begin{split}
 \tau_{\mathrm{full}}
 &=\frac{D}{2(1-DM)}
   =\frac{g^2}{4}(1+O(\varepsilon_X))
   =\frac{4\pi^2}{L^2}(1+O(\varepsilon_X)),\\
 \tau_3
 &=\frac{D}{2(1-Db_3)}
   =\frac{g^2}{6}(1+O(\eta_X))
   =\frac{8\pi^2}{3L^2}(1+O(\eta_X)).
 \end{split}
 \label{ct:horizons}
\end{equation}
In particular, for every fixed $\Delta>0$, even the exact-data score
is negative at all sufficiently high such pairs.
\end{theorem}

\begin{proof}
Apply Lemma~\ref{ct:head-tail} with $R=1$. The replacement
$d=(4\pi/L)(1+O(X^{-1}))$ is absorbed in $\eta_X$.
The complete tail error is $O(L^{-1})$, giving the stated error for
$DM$. The neighboring gaps are $d(1+O(\eta_X))$, so each term
of $b_3$ is $(2d^2)^{-1}(1+O(\eta_X))$; this gives
$Db_3=1/4+O(\eta_X)$.

For fixed $\Delta\ge0$, the score
$R(d_0,b_0;\Delta)=d_0(1+2b_0\Delta)-2\Delta$ is nondecreasing
in both nonnegative lower inputs. Exact $D,M$ therefore maximize it
in this information class. Eventually $DM<1$ and $Db_3<1$;
solving the strict score inequalities and using
\eqref{ct:initial-asymptotics} proves \eqref{ct:horizons}. More
explicitly, at fixed positive $\Delta$,
\begin{equation}
 R(D,M;\Delta)=D-\Delta+O(\Delta\varepsilon_X)\longrightarrow-\Delta,
 \qquad
 R(D,b_3;\Delta)\longrightarrow-3\Delta/2.
 \label{ct:score-limits}
\end{equation}
The errors are uniform over the stated slices.
\end{proof}

The full initial mass improves the asymptotic allowance by a factor
$3/2$ over two neighbors, and by a factor $2$ over the gap-only
$g^2/8$. Both improved durations still vanish as $L^{-2}$. Thus a
nonempty finite block made entirely of sufficiently high eligible pairs
is bad for the fixed-duration sufficient score with probability one
under the finite block sampling of Section~\ref{sec:statistics}.
This concerns the score event, not an actual collision event.

The limitation is confined to the one-step Riccati lower envelope using
initial $D,M$. It excludes neither new time-dependent information nor
reinitialization, divided steps or direct source control. For an exact
check, the field
$H(t,z)=\e^{\pi^2t/g^2}\cos(\pi z/g)$ has stationary zeros,
$D=g^2/4$, $M=2/g^2$ and $D_u=0$, while that same comparison stops
at $g^2/4$. The discarded forcing balances $2M^2$ exactly. The field
has linear zero count and is not the canonical theta source; its purpose
is to distinguish an allowance from a dynamical lifetime.

\subsection{A growing head pays the joint comparison}

Return to the fixed common labels on $I$, let $u=t_*-t$, and suppose
the complete divisor at $t_*$ is real and simple. For this subsection
apply positive omission only on its globally real simple stopped part.
Set
\begin{equation}
 R=L^2,\qquad
 J=\left\lfloor\frac{RL}{64\pi}\right\rfloor,\qquad
 E_X=LX^{-c_*t_0}+\frac{L^4}{X}.
 \label{ct:growing-head}
\end{equation}

\begin{theorem}[Canonical acquisition of the joint barrier]
\label{ct:joint-acquisition}
There are constants $C_K,C_M$ such that, uniformly on that stopped
regime,
\begin{equation}
 K\ge1-C_K(E_X+L^{-9}),\qquad
 |V|\le C_M(E_X+L^{-3}).
 \label{ct:acquired-inputs}
\end{equation}
Consequently, with
\begin{equation}
 \eta=\sqrt{C_K(E_X+L^{-9})},\qquad
 v_X=\max\{C_M(E_X+L^{-3}),\eta\},
 \label{ct:barrier-parameter}
\end{equation}
the joint comparison gives
\begin{equation}
 D(u)\ge D(0)-2v_Xu,\qquad
 v_X=O(L^{-3}+\sqrt{E_X})=O(L^{-3}).
 \label{ct:joint-output}
\end{equation}
Every fixed duration within the stipulated slab satisfies the strict
margin eventually, since
\begin{equation}
 \frac{2v_X\Delta}{D(0)}=O(\Delta/L)\longrightarrow0.
 \label{ct:relative-margin}
\end{equation}
\end{theorem}

\begin{proof}
The fixed head contains labels $n-J,\ldots,n+J+1$. All lie within
distance $R$ of $a=x_n(t)$ for every $t\in I$, after increasing $X$.
Indeed $G_x\ge L/(8\pi)$ on the radius-$2R$ window, so a distance
$R$ allows at least $RL/(8\pi)$ clock indices. This is much larger
than $J+1$, and the source displacements are $o(1)$. The slack pays
both the endpoint asymmetry and the moving spatial buffers. Also
$J=O(L^3)=o(n)$, so the leftmost label still meets the source index
onset. The same finite labels, rather than a moving cutoff, are
therefore available throughout the slab.

Lemma~\ref{ct:head-tail} supplies
$(x_{n+k}-a)/d=k+e_k$ with $|e_k|\le CE_X$ for all these labels.
It follows that $g\ge(1-e)d$ and the two endpoint-distance
inequalities \eqref{jt:head-distances} hold with $e=C'E_X$:
subtracting two normalized positions costs at most $2CE_X$, and
each reference integer distance is at least one. Eventually $e<1/3$.
Proposition~\ref{jt:robust-head}, together with $J\asymp L^3$,
gives the first inequality in \eqref{ct:acquired-inputs}. Its omitted
forcing terms are positive on the stopped regime. No upper tail
subtraction or independent finite-divisor evolution has been used.

For $M$, use the complete statement \eqref{ct:complete-mass} at
$R=L^2$. It gives $2DM=1+O(E_X+L^{-3})$, including the unbounded
tail payment \eqref{ct:tail-payment}, hence the second inequality
in \eqref{ct:acquired-inputs}. Only its initial instance is needed
for Theorem~\ref{jt:barrier}. That theorem gives
\eqref{ct:joint-output}. At fixed $t_0>0$, $E_X$ is smaller than
every fixed negative power of $L$, while
$D(0)\sim4\pi^2/L^2$. These facts prove the last equality in
\eqref{ct:joint-output} and \eqref{ct:relative-margin}.
\end{proof}

When $v_X>0$, the algebraic duration budget $D(0)/(2v_X)$ is at
least a constant times $L$. This is a budget of the comparison; its
inputs have only been established on the fixed slab \eqref{ct:slab}.
It is not a claim of backward evolution for a time of order $L$.
The change from the shrinking allowance in
Theorem~\ref{ct:initial-horizon} is the acquired whole-time forcing
floor, not a more accurate initial mass.

The same-information comparison is already visible:
\eqref{ct:acquired-inputs} also bounds $|V|$ at every time, so
$|D_u|\le2C_M(E_X+L^{-3})$ follows directly, without taking the
square root of a forcing error. More fundamentally, the clock theorem
already proves exterior reality and simplicity independently of the
low divisor. The barrier has a weaker data interface---a full-time
lower $K$ bound plus initial $V$---but no cheaper acquisition or
improved canonical onset is demonstrated here.

\subsection{An analytic source residual and the velocity consequence}

There is a sharper way to retain the common motion of neighboring
zeros. Polymath's Remark~9.4 explicitly describes the spatial-jet
method for a velocity asymptotic and leaves its details to the reader
\cite{Polymath2019DBN}. The proof below supplies the local residual,
Cauchy and denominator estimates needed for this application; the
velocity method is not claimed as new.

Let $q=q_n(t)$, $L_q=\log(q/(4\pi))$, and define
\begin{equation}
 \lambda_q=L_q/4,\qquad \kappa_0=\pi/8,\qquad
 \epsilon_q=q^{-c_*t_0}+L_q/q.
 \label{ct:jet-scale}
\end{equation}
Write $\mathfrak M_t$ for the nonvanishing carrier denoted $M_t$ in
\cite{Polymath2019DBN}; it is not the exterior interaction $M$.
Set $A_q=\mathfrak M_t((1+\ii q)/2)\ne0$.

\begin{lemma}[Holomorphic local residual]
\label{ct:analytic-residual}
For a sufficiently small fixed $r>0$, uniformly for $t\in I$ and
all sufficiently high clock labels,
\begin{equation}
 H(t,q+w)=A_q\e^{-\kappa_0w}
          \{2\ii\sin(\lambda_qw)+E_{n,t}(w)\},
 \quad |w|\le r/L_q,
 \quad \sup|E_{n,t}|\le C\epsilon_q,
 \label{ct:residual}
\end{equation}
where $E_{n,t}$ is an actual holomorphic function on the disk.
\end{lemma}

\begin{proof}
Keep both analytic carrier terms in \cite[\S9, (102)]{Polymath2019DBN}.
For $w=(x-q)+\ii y$, their exponents have leading parts
\[
 (-\ii w)(\lambda_q-\ii\kappa_0)
       =-\ii\lambda_qw-\kappa_0w,
 \qquad
 (\ii w)(\lambda_q+\ii\kappa_0)
       =\ii\lambda_qw-\kappa_0w.
\]
The corresponding multipliers are $\overline{A_q}$ and $A_q$.
The relative errors in those two terms are $O(q^{-c_*t_0})$ on the
disk, with the fixed small source exponent chosen above.
Lemma~9.2 and (103) of that source give
$\overline{A_q}/A_q=-1+O(L_q/q)$ at the clock point.
The two exponentials have uniformly bounded modulus on
$|w|\le r/L_q$. Divide the source approximation by
$A_q\e^{-\kappa_0w}$: the leading terms combine to
$\e^{\ii\lambda_qw}-\e^{-\ii\lambda_qw}$ and the remaining
terms are bounded by $C\epsilon_q$.

Define $E_{n,t}$ by the exact left side of \eqref{ct:residual}
minus $2\ii\sin(\lambda_qw)$. The quotient uses a nonvanishing
entire factor, so this definition proves holomorphy, and the source
estimate proves its uniform bound. No derivative of a symbol $O(\cdot)$
is taken. The leading frequency is deliberately $L_q/4$; the smaller
time-dependent phase refinement is included in the stated error.
The weaker display (104) in the source, with an allowance involving
$|\Im w|\log q$, would not give a vanishing residual on this disk
and is not used for the Cauchy step.
\end{proof}

\begin{proposition}[Source jets and the complete gap defect]
\label{ct:velocity}
For the common high labels on $I$, put
$\epsilon_X=X^{-c_*t_0}+L/X$. Without an assumption that the low
divisor is real,
\begin{equation}
 \begin{split}
 |x_n'(t)+\pi/4|&\le C L\epsilon_X,\\
 |g_t|&\le C L\epsilon_X,\qquad
 |D_t|\le C\epsilon_X,\qquad
 |1-2DM|\le C\epsilon_X.
 \end{split}
 \label{ct:velocity-bounds}
\end{equation}
Here $M$ is the complete exterior sum for the adjacent high pair.
In particular, for $s,t\in I$,
\begin{equation}
 |g(t)-g(s)|\le C L\epsilon_X|t-s|,
 \qquad |D(t)-D(s)|\le C\epsilon_X|t-s|.
 \label{ct:integrated-jets}
\end{equation}
\end{proposition}

\begin{proof}
Source localization puts the actual real root
$w_n=x_n(t)-q_n(t)$ in $|w_n|\le r/(4L_q)$ for sufficiently large
$q$, uniformly in $t$. Equation~\eqref{ct:residual} at that root
implies $\sin(\lambda_qw_n)=O(\epsilon_q)$. On this fixed small
real interval the sine is bounded below in modulus by a constant
times its argument, so
\begin{equation}
 |w_n|\le C\epsilon_q/L_q.
 \label{ct:sharpened-position}
\end{equation}
Cauchy's formula on the larger disk gives, in its inner half,
\[
 |E'_{n,t}|\le CL_q\epsilon_q,\qquad
 |E''_{n,t}|\le CL_q^2\epsilon_q.
\]
Put $f(w)=2\ii\sin(\lambda_qw)+E_{n,t}(w)$. At the root,
$\lambda_qw_n=O(\epsilon_q)$, and therefore
\begin{equation}
 f'(w_n)=2\ii\lambda_q(1+O(\epsilon_q)),\qquad
 f''(w_n)=O(L_q^2\epsilon_q),\qquad
 |f'(w_n)|\ge c_1L_q.
 \label{ct:root-jets}
\end{equation}
The second bound uses the small sine at the root and the paid
second derivative of the residual. Spatially differentiating the
exact representation \eqref{ct:residual}, with $t,q,A_q$ fixed,
and then setting $f(w_n)=0$ gives
\[
 \frac{H_{zz}(t,x_n)}{H_z(t,x_n)}
 =\frac{f''(w_n)}{f'(w_n)}-2\kappa_0
 =-\frac\pi4+O(L_q\epsilon_q).
\]
Implicit differentiation of $H(t,x_n(t))=0$ and $H_t=-H_{zz}$
gives the exact physical velocity $x_n'=H_{zz}/H_z$.
This proves its asserted bound. The leading velocity $-\pi/4$
cancels for neighboring labels, so $g_t=O(L\epsilon_X)$;
since $g=O(L^{-1})$, $D_t=gg_t/2=O(\epsilon_X)$.

For completeness, this derivative estimate transfers to the complete
$M$, not only to a local head. The high pair is strictly isolated;
apply the finite-product pair identity at a reference time and pass
through the cleared pair contour as in Lemma~\ref{lem:slice-passage}.
Equivalently, subtract the two finite root-velocity sums before their
complete opposite/conjugate-group limit. The exterior difference is
\[
 g_t=\frac4g-2g\sum_{\rho\ne a,b}
                       \frac1{(\rho-a)(\rho-b)}.
\]
Every surviving tail term is inverse-square and hence absolutely
summable by the complete strip and unit count. The nearby denominators
are nonzero by strict isolation, including when lower zeros are
nonreal. Thus $D_t=2-4DM=2V$ holds with the complete real sum,
proving the last bound in \eqref{ct:velocity-bounds}. No separate
nonabsolute first-power sum is estimated. Integration on the common
label branches proves \eqref{ct:integrated-jets}.
\end{proof}

This proof uses complex spatial neighborhoods and the exact heat
equation, not a complex-time source estimate or a time derivative of
the localization error. Since $\epsilon_X L^2\to0$, its squared
half-gap allowance pays every fixed duration in $I$ at sufficiently
high labels. The corresponding allowance for the actual gap has the
additional factor $L$ displayed in \eqref{ct:integrated-jets}.

\subsection{The matched direct clock comparison}

The same source information also encloses the gaps at each time.
Let $h_n(t)=q_{n+1}(t)-q_n(t)$. From
\eqref{ct:sharpened-position},
\begin{equation}
 |g_n(t)-h_n(t)|\le C\epsilon_X/L.
 \label{ct:direct-enclosure}
\end{equation}
This is a common-time enclosure, not merely an initial-value estimate.
Write $v(q,t)$ for the explicit clock velocity in
\eqref{ct:clock-velocity}, and $L_q=\log(q/(4\pi))$. Differentiating
this explicit function gives
\[
 v_q(q,t)=-\frac{\pi^2t(L_q+1)}{16q^2L_q^2}
          \left(1+\frac{\pi t}{4qL_q}\right)^{-2}.
\]
Since $h_n=O(L^{-1})$, the mean-value theorem implies
$|h_n'(t)|\le C/(X^2L^2)$. Combining two endpoint enclosures
\eqref{ct:direct-enclosure} with this bound yields
\begin{equation}
 \begin{split}
 |g_n(t)-g_n(s)|
 &\le C\epsilon_X/L+C|t-s|/(X^2L^2),\\
 |D_n(t)-D_n(s)|
 &\le C\epsilon_X/L^2+C|t-s|/(X^2L^3).
 \end{split}
 \label{ct:matched-direct}
\end{equation}
For the second line, use $D=g^2/4$ and $g(t)+g(s)=O(L^{-1})$.
Both endpoint source errors are included in the constants. For a fixed
positive time separation this direct bound pays roughly $L^2$ less
source error than the integrated squared half-gap bound. For very short
times, \eqref{ct:integrated-jets} can be better because its allowance
vanishes with $|t-s|$. The minimum of the two bounds is legitimate;
neither dominates in every time regime.

These comparisons locate the gain and its boundary. The original
initial-$M$ envelope loses a cancellation; the acquired $K$ floor
retains enough of it to pay a fixed slab, and the analytic source
residual estimates it still more sharply. All source acquisitions
remain in the exterior already covered by the published theorem.
They give no numerical cutoff, no estimate uniform down to $t=0$,
no new compact gap law, and no improved bound on $\Lambda$. Extending
the useful domain requires new source information of the type considered
in Section~\ref{sec:compact}, rather than repeated evaluation of this
already protected high exterior.
